\section{PNEW-016: a dimension bound for the surviving tensor waves}
\label{sec:tensor-fixed-dimension-gate}

Put $V_n=(\mathbb C^n)^{\otimes3}$, $G_n=SL_n^3$, and
$X_n=V_n//G_n$. The grading on its coordinate ring is ordinary
polynomial degree divided by $n$. On $GL_n^3$ it is induced by
$\chi(A,B,C)=\det A\det B\det C$. Write $X_n^{\mu_d}$ for the
fixed locus of its grading subgroup of order $d$.

For $d\ge2$, define the integer
\[
 r(d)=\min\{r\ge0:3^r\ge d^2\}.
                                                        \tag{PNEW-I1}
\]
The following bound concerns the actual tensor quotient, including
points with positive-dimensional special-linear stabilizer.

\begin{theorem}[PNEW-016a]
If $X_n^{\mu_d}$ contains a nonzero point, then $r(d)\le3n-2$, and
\[
 \dim X_n^{\mu_d}
 \le \left(n-\frac{r(d)+2}{3}\right)^3+3n^2.
                                                        \tag{PNEW-I2}
\]
In particular, for $M=2^{n-1}$ and $n\ge2000$,
\[
 \dim X_n^{\mu_M}\le\frac{n^3}{5}.
                                                        \tag{PNEW-I3}
\]
The leading coefficient supplied by the sharper bound (PNEW-I2) is
$\bigl(1-2/(3\log_2 3)\bigr)^3$, approximately $0.19449$.
\end{theorem}

\begin{proof}
Let $x\ne0$ be fixed by a primitive $d$th grading root $\xi$ and
choose the closed $G_n$-orbit $G_nv$ above it. Choose $l_0\in GL_n^3$
with $\chi(l_0)=\xi$. Since $l_0v$ and $v$ have the same quotient
and closed special-linear orbits, $l_0v=gv$ for some $g\in G_n$.
Thus $l=g^{-1}l_0$ fixes $v$ and has character $\xi$. Its semisimple
Jordan part also fixes $v$ and has the same character: the unipotent
part has determinant one. Replace $l$ by this semisimple part and
conjugate it to diagonal form in each factor.

Let its three lists of distinct eigenvalues have lengths
$m_1,m_2,m_3$, with multiplicities $d_{a,i}\ge1$ summing to $n$
in each factor. Let $E$ be the supported block triples of $v$.
Their incidence columns lie in $\mathbb Z^{m_1+m_2+m_3}$ and each
has three entries equal to one. Choose a positive-degree homogeneous
special-linear invariant nonzero at $v$. At least one monomial in
it is nonzero at $v$. Each of its monomials has weight zero for the
three special-linear diagonal tori, so its coordinate exponents have
equal margins within each factor. If its degree is $ns$, those
margins are $s>0$. Dividing the exponents by $s$ gives rational
nonnegative margin-one masses on supported coordinates. Aggregating
over the eigenvalue blocks therefore writes
\[
       w=(d_{1,1},\ldots,d_{1,m_1};d_{2,1},\ldots;
                           d_{3,1},\ldots)
          =\sum_{e\in E}q_e a_e,
             \qquad q_e\ge0.                         \tag{PNEW-I4}
\]
In particular $w$ is torsion modulo the incidence-column lattice $L$
in its rational span. This necessity uses only a nonzero invariant
monomial, not a sufficient-polystability criterion for supports.

Let $r=\operatorname{rank}L$ and choose $r$ independent original
incidence columns, with column matrix $B$. Their lattice $L'$
is contained in $L$, so the torsion order of $w$ modulo $L$ is
at most $[\operatorname{sat}(L'):L']$. The latter index is the
greatest common divisor of the nonzero $r$ by $r$ minors of $B$.
It is at most $\sqrt{\det(B^tB)}$, by Cauchy--Binet, and then at
most $3^{r/2}$, by Hadamard, since each column has norm $\sqrt3$.
Evaluation at the diagonal eigenvalues kills every supported
incidence column, whereas evaluation at $w$ is $\chi(l)=\xi$.
Consequently that torsion order is divisible by $d$, and
\[
             d\le3^{r/2},\qquad
             r(d)\le r\le m_1+m_2+m_3-2.             \tag{PNEW-I5}
\]
The last inequality uses the two independent relations equating
the sums of the rows in the three factors.

Now enlarge $E$ to the complete fixed-block support $E^*$ of this
semisimple element. Given two block indices, at most one third
index lies in $E^*$, because the eigenvalues within a factor are
distinct. Write $d_{1,i}=1+\alpha_i$, $d_{2,j}=1+\beta_j$,
$d_{3,k}=1+\gamma_k$, with all entries nonnegative. For each edge,
the expansion of $(1+\alpha_i)(1+\beta_j)(1+\gamma_k)$ apart from
its cubic term is bounded by the sum of its three two-factor
products. Pair uniqueness then gives
\[
\begin{split}
 \dim V_n^l
 &=\sum_{(i,j,k)\in E^*}
        (1+\alpha_i)(1+\beta_j)(1+\gamma_k)\\
 &\le \sum_{(i,j,k)\in E^*}\alpha_i\beta_j\gamma_k+3n^2\\
 &\le (n-m_1)(n-m_2)(n-m_3)+3n^2\\
 &\le\left(n-\frac{r(d)+2}{3}\right)^3+3n^2.
\end{split}                                             \tag{PNEW-I6}
\]
For the final step, use (PNEW-I5) and arithmetic--geometric mean
on the three nonnegative numbers $n-m_a$.

There are only finitely many ordered multiplicity lists and
fixed-block supports. Thus the diagonal fixed spaces appearing
above belong to a finite list of coordinate-support subspaces.
Conjugation by $GL_n^3$ changes the special-linear quotient only
by grading. The image of any such linear subspace is a cone,
so this grading translation does not enlarge its image. Hence
$X_n^{\mu_d}$ is covered by the origin and finitely many images
of the subspaces just bounded. The dimension of each image is at
most its source dimension, proving (PNEW-I2). This argument does
not require a uniform bound on the order of the lift itself.

For the last assertion, $3^5=243<256=4^4$ implies
$r(M)>5(n-1)/4$. Therefore
\[
 \dim X_n^{\mu_M}
 \le\left(\frac{7n-3}{12}\right)^3+3n^2
 \le\frac{343n^3}{1728}+3n^2\le\frac{n^3}{5}
                                                        \tag{PNEW-I7}
\]
when $n\ge2000$; the last inequality follows from
$13n\ge25920$.
\end{proof}

\subsection*{An actual covariant difference of smaller degree}

Retain $k_n=b_n+3c_n$ and the actual hook wave of PNEW-015.
Fix any common period $P=2^vQ$, with $Q$ odd, of these three
global quasipolynomials. Necessarily $v\ge n-1$. Put
\[
\begin{gathered}
 h=(M/2)Q,\qquad
 \Delta_{n,j}(X)=c_{n,j+h}(X)-c_{n,j}(X),\\
 U_{\Delta}^+(n;P)=\max_{j\bmod P}\sum_{a\ge0}
                   \max\{[X^a]\Delta_{n,j}(X),0\}.
\end{gathered}
                                                        \tag{PNEW-I8}
\]
The period is fixed before applying the next statement. No independence
of this particular difference under a change of common period is
asserted.

\begin{lemma}[PNEW-016b]
For $n\ge2000$, every nonzero $\Delta_{n,j}$ has degree at most
$n^3/5-1$. In particular the common polynomial part at the root
one makes no contribution to (PNEW-I8).
\end{lemma}
\begin{proof}
Write the ordinary quasipolynomial by its root-of-unity partial
fractions. A pole at $\zeta$ of order $e$ contributes a polynomial
of degree at most $e-1$ multiplied by the residue phase
$\zeta^{-j}$. Shifting $j$ by $h$ kills every root whose order has
two-part smaller than $M$, since its odd part divides $Q$.
Every surviving root has order divisible by $M$.

For completeness, the pole order of a graded coordinate ring at
$\zeta$ is at most the dimension of its $\zeta$-fixed locus.
Take homogeneous algebra generators and the Koszul complex on
those on which $\zeta$ acts nontrivially. Its Euler identity
multiplies the Hilbert series by a factor nonzero at $\zeta$.
Its homology modules are supported on the fixed locus and are
modules over the quotient generated by the remaining, fixed
generators. Decomposing their shifts into grading residues shows
that their pole orders at $\zeta$ are at most their dimensions,
exactly as at the root one. This proves the claim, including for
singular fixed loci. It is the pole-order form of the integral
Koszul argument already used in K-027.

For every surviving tensor root the fixed locus lies in
$X_n^{\mu_M}$, so (PNEW-I3) bounds its pole order by $n^3/5$.
K-024 says that the surviving hook roots are precisely the
primitive-$M$ simple poles. Since $c_n=(k_n-b_n)/3$, the same
degree bound holds for their difference (PNEW-I8).
\end{proof}

\begin{theorem}[PNEW-016c]
For $n\ge79004$ and any fixed common period $P$ as above, either
\[
        p(k_n)>2^{n/16},\qquad s(k_n)\ge2^{n/16},
                                                        \tag{PNEW-I9}
\]
both hold, or
\[
       U_{\Delta}^+(n;P)\ge 2^{(17/240)n^4-6n^3-4}.              \tag{PNEW-I10}
\]
Thus a bound $U_{\Delta}^+(n;P)\le2^{C n^3\log_2(n+2)}$ for some common
period and a fixed $C\ge1$ implies both conclusions in
(PNEW-I9) whenever
\[
       C\log_2(n+2)+8\le17n/240.                     \tag{PNEW-I11}
\]
No bound of this kind on $U_{\Delta}^+(n;P)$ is proved here.
\end{theorem}
\begin{proof}
Let $A_n\ge2^{n^4/12-5n^3}$ be the selected hook-wave amplitude
of PNEW-015. Choose an integer $x$ in its nine-argument phase
window beginning at $\lceil2^{n/16}\rceil$, for which
$2W_n^b(x)\ge A_n/2$. Then $x\le2^{n/16+1}$, and the actual
tensor ambient bound gives
\[
 k_n(x)\le\binom{nx+n^3-1}{n^3-1}
       \le2^{n^3(n/16+2+\log_2n)}<A_n/4.             \tag{PNEW-I12}
\]
Indeed $x\ge n^2$ in this range,
$2+\log_2n\le n/96$, and $n^4/96>5n^3+2$.

Put $j=x\bmod P$. The exact hook phase identity (PNEW-H12)
now gives, for the actual tensor constituents,
\[
 k_{n,j+h}(x)=k_n(x)-2W_n^b(x)+3\Delta_{n,j}(x).
                                                        \tag{PNEW-I13}
\]
If either bound in (PNEW-I9) fails, the corresponding strict-index
definition makes the left side nonnegative. Consequently
$\Delta_{n,j}(x)>A_n/12$. Discarding negative coefficients in
this polynomial, and using PNEW-016b, gives
\[
                 \Delta_{n,j}(x)\le U_{\Delta}^+(n;P)x^{n^3/5}.
\]
It follows that
\[
\begin{split}
 \log_2 U_{\Delta}^+(n;P)
 &>n^4/12-5n^3-\log_2 12-(n^3/5)(n/16+1)\\
 &\ge (17/240)n^4-6n^3-4.
\end{split}
\]
This proves the dichotomy. Condition (PNEW-I11) makes the proposed
upper bound strictly smaller than its second alternative, and
therefore yields (PNEW-I9).
\end{proof}

\paragraph{A Fourier limitation of the proposed coefficient promise.}
For a primitive $M$th root $\zeta$, define its actual tensor
constituent Fourier polynomial by
\[
 p_{k,\zeta}(X)=\frac1P\sum_{j\bmod P}k_{n,j}(X)\zeta^j.
\]
The same notation applies to $b,c,\Delta$. Write
$C_{\rm hook}^{\rm gr}(\zeta)$ for the moving coefficient
indexed by its grading root: it equals $C_{\rm hook}(\eta)$
of PNEW-015 when $\eta^S=\zeta$ and
$S=\sum_{i=0}^{n-2}(-2)^i$. Since $\zeta^h=-1$, we have
\[
 p_{\Delta,\zeta}(X)
   =-\frac23\left(p_{k,\zeta}(X)
                       -\frac{C_{\rm hook}^{\rm gr}(\zeta)}M\right).
                                                        \tag{PNEW-I14}
\]
For each exponent $a$, the real numbers
$d_{j,a}=[X^a]\Delta_{n,j}$ sum to zero over all residues.
Consequently
\[
 \left|\frac1P\sum_jd_{j,a}\zeta^j\right|
 \le\frac1P\sum_j|d_{j,a}|
 =\frac2P\sum_j\max(d_{j,a},0)
 \le2U_\Delta^+(n;P).
\]
In particular the proposed promise would imply
\[
 \left|[X^0]p_{k,\zeta}(X)
                       -\frac{C_{\rm hook}^{\rm gr}(\zeta)}M\right|
       \le3U_\Delta^+(n;P).                            \tag{PNEW-I15}
\]
If the actual tensor series has no pole at the selected primitive
$M$th root, its Fourier polynomial is zero. The already proved
amplitude bound therefore forces
\[
 U_\Delta^+(n;P)
      \ge\frac{|C_{\rm hook}^{\rm gr}(\zeta_{\rm selected})|}{3M}
       =A_n/3
        \ge\frac13\,2^{n^4/12-5n^3}.                  \tag{PNEW-I16}
\]
Precisely, $\zeta_{\rm selected}=\exp(2\pi i9^{-1}/M)$ is
the grading root; $A_n=|C_{\rm hook}^{\rm gr}(\zeta_{\rm selected})|/M$
is the single-root magnitude defined in PNEW-015, not twice it.
Thus a sufficiently small subquartic upper bound on this norm
would also prove exact-$M$ noncancellation and force a large constant
Fourier coefficient of the actual tensor series. It is a strong
spectral-local estimate, not a consequence of the smaller degree.
No such estimate or exact-$M$ conclusion is inferred here.

\paragraph{Scope and provenance.}
The incidence-lattice and Hadamard argument is the standard mechanism
behind K-009, applied here to the compressed eigenspace multiplicities
to control fixed-space dimension. The margin necessity is proved
directly from an invariant monomial above; the closed-orbit quotient
property and integral Koszul method are the previously checked GIT
inputs. The proof supplies the new combination
and all of its inequalities. This is an actual tensor fixed-locus and
constituent-degree bound, followed by a conditional coefficient
estimate. It does not compute the remaining coefficient mass or
settle the strict tensor indices. In particular, the smaller degree
does not by itself bound the size of polynomial coefficients.

\paragraph{Bounded primary-literature comparison.}
The determinant-product stabilizer period and the support-margin
framework are established tensor invariant theory. A directly checked
primary reference is B\"urgisser, P., \& Ikenmeyer, C. (2017).
Fundamental invariants of orbit closures.
\emph{Journal of Algebra, 477}, 390--434.
\url{https://doi.org/10.1016/j.jalgebra.2016.12.035}.
The exact reading scope here was the character and period definition
in Section~4.1, equations~(4.1)--(4.2) and Definition~4.1, printed
pp.~18--19 of the author version, and the full support-margin setup,
Proposition~4.8 and its proof in Section~4.2, printed pp.~21--22:
\url{https://arxiv.org/pdf/1511.02927}.
These passages supply the inspected framework; neither their
sufficient-polystability assertion nor their printed averaging
argument is an input to our proof. In particular a probability
distribution there need not be positive on every supported edge.
They are not cited as a source for the compressed dimension estimate
(PNEW-I2), whose margin necessity was proved independently above.

For comparison with known large incidence denominators, the previous
primary audit read Franks, W.~C., \& Reichenbach, P. (2021).
Barriers for recent methods in geodesic optimization. In V.~Kabanets
(Ed.), \emph{36th Computational Complexity Conference}
(LIPIcs, Vol.~200, Article~13, pp.~13:1--13:54).
Schloss Dagstuhl--Leibniz-Zentrum f\"ur Informatik.
\url{https://doi.org/10.4230/LIPIcs.CCC.2021.13}.
The complete Lemma~2.5 and proof, printed p.~13:15, give a known
exponential-denominator construction and explicitly credit earlier
axial-assignment literature. This is background for why large orders
must be retained, not a dimension estimate for all fixed components.
The earlier audit records the exact original-source reading limits.

Related general coefficient-grade control appears in Bruns, W., \&
Ichim, B. (2007). On the coefficients of Hilbert quasipolynomials.
\emph{Proceedings of the American Mathematical Society, 135}(5),
1305--1308, Theorem~2 and Lemmas~3--4, pp.~1306--1307:
\url{https://www.home.uni-osnabrueck.de/wbruns/brunsw/pdf-article/Quasipol.published.pdf}.
All four pages of this primary PDF were read. Their ideal generated
in degrees coprime to the period controls which coefficient functions
are constant. This is related background for the support/degree
mechanism, not a source for the tensor estimate (PNEW-I2) or an
upper bound on constituent coefficient sizes.

The determinant inequality used in (PNEW-I5) is Hadamard's classical
Gram determinant inequality: Hadamard, J. (1893).
R\'esolution d'une question relative aux d\'eterminants.
\emph{Bulletin des Sciences Math\'ematiques, 17}, 240--246
(2nd series), Sections~2--4, pp.~240--242.
Those sections were checked in the university-hosted English
translation at
\url{https://maths.nuigalway.ie/~padraig/Docs/HadamardTranslation.pdf};
no fresh reading of the original French scan is claimed.
This was a bounded comparison of the indicated sources and methods,
not an exhaustive priority search. No historical priority claim is
made for (PNEW-I2)--(PNEW-I3).
